\section{Stratification Methodology}
\label{app:stratification}

To ensure representative coverage in the development and holdout splits, we employed stratified sampling across three complementary dimensions: semantic category, task complexity, and difficulty level.

\subsection{Stratification Dimensions}

\paragraph{Dimension 1: Semantic Category}
Prompts are classified into three broad categories using keyword-based heuristics:

\begin{itemize}
    \item \textbf{STEM}: Mathematical reasoning, scientific queries
    \begin{itemize}
        \item Keywords: \textit{integral, derivative, theorem, proof, math, calculus, equation}
    \end{itemize}
    
    \item \textbf{CODE}: Programming tasks, debugging, technical implementation
    \begin{itemize}
        \item Keywords: \textit{code, function, class, debug, python, javascript, rust, ```}
    \end{itemize}
    
    \item \textbf{GENERAL}: All other prompts (conversational, creative, knowledge)
    \begin{itemize}
        \item Default category when no category-specific keywords match
    \end{itemize}
\end{itemize}

\paragraph{Dimension 2: Task Complexity}
Complexity is estimated using structural features of the prompt:

\begin{itemize}
    \item Length-based signals (long prompts $>500$ chars, very long $>2000$ chars)
    \item Presence of code blocks (``` markers)
    \item Reasoning indicators (\textit{step-by-step, explain, why, how})
    \item Constraint keywords (\textit{without, avoid, only, constraint})
\end{itemize}

A composite score $\in [0, 1]$ is computed, yielding three complexity levels:
\begin{itemize}
    \item \textbf{Low}: score $< 0.3$
    \item \textbf{Medium}: $0.3 \leq$ score $< 0.7$
    \item \textbf{High}: score $\geq 0.7$
\end{itemize}

\paragraph{Dimension 3: Difficulty Level}
Difficulty is inferred from the oracle reward distribution for each prompt:

\begin{itemize}
    \item \textbf{Easy}: High mean reward ($\mu > 0.8$), low variance ($\sigma^2 < 0.05$)
    \item \textbf{Hard}: Low mean reward ($\mu < 0.8$), low variance ($\sigma^2 < 0.05$)
    \item \textbf{Contentious}: High variance ($\sigma^2 \geq 0.05$), indicating disagreement between models
\end{itemize}

\subsection{Stratification Procedure}

Each prompt is assigned a stratification key combining all three dimensions:
\begin{equation}
\text{key} = \text{Category}_{\text{Complexity}}_{\text{Difficulty}}
\end{equation}

For example: \texttt{STEM\_High\_Contentious}, \texttt{CODE\_Med\_Easy}, \texttt{GENERAL\_Low\_Hard}.

\paragraph{Sparse Strata Handling.}
Strata with fewer than 2 members are merged into fallback categories to ensure valid stratified splitting (sklearn's \texttt{train\_test\_split} requires $\geq 2$ samples per stratum):
\begin{itemize}
    \item First-level fallback: \texttt{Category\_Mixed\_Mixed}
    \item Second-level fallback: \texttt{GENERAL\_Mixed\_Mixed}
\end{itemize}

\paragraph{Split Generation.}
Using stratified sampling (60\% development, 40\% holdout with random seed 42), we generate splits that maintain proportional representation of each stratum across both sets. This ensures:
\begin{itemize}
    \item Similar distributions of semantic categories
    \item Similar distributions of task complexity
    \item Similar distributions of difficulty levels
\end{itemize}

\subsection{Verification}

The stratification is verified via:
\begin{enumerate}
    \item \textbf{Disjointness check}: Zero overlap between dev and holdout (243 duplicates removed pre-split)
    \item \textbf{Stratum balance}: Each stratum represented proportionally in both splits
    \item \textbf{Distribution similarity}: Chi-square tests on category distributions (though see limitations below)
\end{enumerate}

\subsection{Limitations}

\paragraph{Heuristic-Based Categories.}
The semantic categories (STEM, CODE, GENERAL) are determined by keyword matching, not semantic understanding. While this provides a coarse but reproducible categorization, it may misclassify edge cases (e.g., "Explain how sorting algorithms work" could be CODE or GENERAL depending on phrasing).

\paragraph{No Ground Truth.}
The complexity and difficulty dimensions are heuristically computed, not validated against human annotations. They provide structural diversity in sampling but should not be interpreted as validated task taxonomies.

\paragraph{Different From Archived Analysis.}
Earlier versions of the dataset analysis used a different 5-category scheme (Coding, Math/Logic, Creative, Knowledge, Conversational). The current implementation uses the simpler 3-category scheme (STEM, CODE, GENERAL) shown here. The categories are used solely for stratification, not for any downstream experimental analysis or claims about category-specific performance.

\subsection{Rationale}

\paragraph{Why Stratify?}
Stratified sampling ensures that rare prompt types (e.g., STEM tasks at 10-15\% of data) are proportionally represented in both development and holdout sets. Without stratification, random sampling could produce imbalanced splits where one set has disproportionately more of a particular type.

\paragraph{Why Three Dimensions?}
Single-dimension stratification (e.g., complexity only) could produce splits balanced on complexity but imbalanced on category. Multi-dimensional stratification provides better coverage across the joint distribution of prompt characteristics.

\paragraph{Does This Affect Results?}
No. The stratification is used solely for data splitting to ensure representative samples. The main experimental results (bandit learning, routing performance) do not depend on category labels and are evaluated uniformly across all prompts regardless of their stratification key.
